% 第 9 章　能量：变化的记账法（完整样章）
\chapter{能量：变化的记账法}

\step{反常识开场}

\begin{mainline}
美国有一位物理教授做过一个著名的课堂演示。他把一个十几千克重的大钢球用粗钢丝绳吊在教室天花板上，做成一个巨大的摆。然后他退到墙边，把钢球拉到自己鼻尖前，贴着皮肤，松手——注意，是松手，不是推。钢球呼啸着荡到教室另一头，又呼啸着荡回来，越回越快，直奔他的脸。全班学生都屏住了呼吸。钢球在离他鼻尖几厘米的地方停住，掉头荡走了。教授纹丝不动。他说：我对物理学有信心，所以我不躲。

这个演示里真正值得失眠的不是教授的胆量，而是钢球的“自律”：它荡回来时，为什么\emph{永远不会}比出发点更高？它不知道那里有鼻子，它只是一块死铁；可它每一次都乖乖停在比出发点略低一点的地方。反过来问：如果松手时偷偷推它一把，它就会越荡越高，真的砸到鼻子。多出来的那一点点“高”，到底是从哪来的？钢球又凭什么把“原来有多高”这件事记得这么牢？

再看一个你随时能观察的版本。公园里荡秋千，不推不蹬，荡幅只会越来越小，最后停在最底下。全世界没有任何一架秋千会自己越荡越高。水往低处流，摆往低处停，好像自然界有一本看不见的账：有些东西可以来回倒手，但总额只减不增，谁也别想凭空多拿。这本账记的是什么？这一章，我们就把账本翻开。
\end{mainline}

\step{先猜一猜}

老规矩，先押注，再读书。请把你的答案写下来。

第一题。两个一模一样的小球，从同一高度由静止滚下：一个沿光滑的直坡，另一个沿光滑的弯弯曲曲、中途还凹下去又拱起来的滑道。不计摩擦和空气阻力。到达底部时，哪个球更快？（a）直坡的快；（b）弯道的快；（c）一样快。

第二题。汽车以 \(40\) 千米每小时行驶，一脚急刹到底，滑行 \(10\) 米停下。同一辆车、同一条路，以 \(80\) 千米每小时急刹，滑行距离大约是多少？（a）\(20\) 米；（b）\(40\) 米；（c）比 \(40\) 米短。

第三题。有人设计了一台机器：一个竖直的轮子，轮缘上挂着一串可以摆动的重锤，设计者说重锤在右边总是“伸得远”、在左边总是“缩得近”，所以右边永远比左边重，轮子会永远转下去，还能顺便带动一台发电机。这台机器能成功吗？如果能，为什么几百年来没有人靠它发电？如果不能，它卡在哪一步？

先写下你的判断。这一章结束时我们回来对账。

\step{动手实验}

\begin{experiment}[钥匙摆：它真的记得高度吗]
\textbf{材料}：一根一米左右的细线（耳机线、缝衣线都行）、一把钥匙或一个螺母、一支笔。

\textbf{步骤一}：把钥匙拴在线的一端，线的另一端用笔画个记号后捏紧，让钥匙竖直下垂。这就是一个摆。

\textbf{步骤二}：把钥匙拉到你的鼻尖旁边，让线绷直，然后\emph{轻轻松手}（千万不要推）。钥匙荡过去再荡回来。注意观察：它回来时最高到达哪里？敢不敢让鼻尖当参照物？（只要你确实没有推，它是安全的——这正是教授演示的家用版。）

\textbf{步骤三}：换一个释放高度再试。无论从哪里释放，钥匙回来时几乎到达\emph{同样}的高度，并且总是略低一点点。多荡几次，逐次降低，最后停在最低点。

\textbf{步骤四}：在线的中途横一根笔杆，让摆锤荡到最低点之后线被笔杆拦住、摆动半径突然变短。你会发现一件惊人的事：钥匙在另一边仍然几乎爬回原来的\emph{高度}，尽管轨迹已经完全不是原来那条弧。

\textbf{记录}：写下你观察到的“它记得的那个量”。是速度吗？最低点速度每次都在变。是路径吗？笔杆一拦路径就变了。可高度它记得清清楚楚。本章要回答的就是：这个世界到底在记一笔什么账。
\end{experiment}

\begin{experiment}[橡皮筋弹射器：两倍的拉伸，几倍的能量]
\textbf{材料}：一根粗橡皮筋、一根小纸棍（把便签纸卷紧）、一把尺子、一本书当发射台。

\textbf{步骤一}：把橡皮筋一端压在书脊下固定，另一端勾住纸棍，沿桌面水平拉开。先用尺量好：拉长 \(2\) 厘米，松手发射，量出纸棍落点到桌边的距离。重复三次取平均。

\textbf{步骤二}：拉长 \(4\) 厘米（整整两倍），同样发射三次取平均。猜猜落点距离会变成几倍？多数人的直觉是“两倍左右”。

\textbf{步骤三}：实测结果通常是三倍上下甚至更多——而落点距离由纸棍出手时的速度决定，速度翻倍才能让射程翻倍。射程变成三倍，说明出手速度大约变成两倍，而纸棍带走的\emph{动能}变成了四倍。

\textbf{记录}：拉伸量只翻了一倍，橡皮筋给出的能量却翻了两倍不止。能量到底按什么规矩存在橡皮筋里？这个平方关系将在“数学翻译器”一节被画成一个三角形。
\end{experiment}

\step{旧模型遇到麻烦}

到上一章为止，我们的全部武器是力：受力图、牛顿三定律、各种具体的力。这套武器回答“此刻”的问题无往不利——此刻受什么力，此刻加速度多大。可是生活里有一大类问题问的不是“此刻”，而是“整段过程累计下来，结果是什么”。拿力的语言去回答这类问题，会立刻露出三处捉襟见肘。

麻烦一：力相同，结果却不讲道理地不同。同样大小的力推同一辆小车，推一秒钟和推两秒钟，速度不同——这好理解，速度跟时间大体成正比地涨。可你再测“这小车还能滑行多远”，就发现不对劲了：速度快一倍的小车，不是滑两倍远，而是滑大约四倍远。如果你的直觉说“力给了小车一份‘冲劲’，冲劲用完了就停”，那么这份冲劲按速度记账和按距离记账，竟然对不上账。旧语言里没有一个量，能同时把“推了多久”和“滑了多远”统一起来。

麻烦二：子弹穿木板。一颗子弹恰好能击穿一块木板，速度降为零。那么两颗同样的子弹、同样的速度，能击穿两块叠起来的木板吗？能——但注意，击穿第一块后速度\emph{不是}降到一半。直觉说“每块木板削掉同样的速度”，可实验和理论都说：木板削掉的是另一种东西，削完那种东西之后，速度只降到原来的七成左右。那种“东西”是什么？力的语言里没有它的位置。

麻烦三：高度为什么被“记住”。钥匙摆换高度、换半径、甚至被笔杆半路拦截，回来时总认原来那条水平线。力每一点都在变（线的拉力方向时刻在变），轨迹一变，力的全程细节就全变了；可按牛顿定律一段段算下去，最后的高度却总是同一个。这像什么？像你沿着七拐八绕的山路开车去邻市，不管走哪条路，油表的读数只取决于起点和终点。山路（力）千变万化，有一个量却纹丝不动——这个量绝不是力本身，它是比力更“大局”的东西。

这三处麻烦有一个共同的形状：力的语言擅长逐瞬间追踪，却不擅长做\emph{全程总结}。而“全程总结”其实有两种切法：把力按\emph{时间}一段段累加起来，得到的东西叫冲量，它管着碰撞，是下一章的主角；把力按\emph{路程}一段段累加起来，得到的正是本章的功。牛顿第二定律本身不偏袒任何一种切法，可大自然偏偏两种都认——于是我们得准备两本账。现在先看按路程记账的这一本：我们需要一种新的记账方式，不问每一瞬间的推力拉力，只问整段过程下来，有什么东西从一个地方搬到了另一个地方、从一种形态换成了另一种形态，而总量分毫不差。这个东西，历史上的物理学家找了将近两百年才找到，期间还给它起过“活力”这样的怪名字。今天它叫能量。

\step{发明新概念}

\begin{mainline}
\textbf{第一个新概念：功——力开出的转账单。}既然要做全程总结，就先定义“一份力在一段路程上总共干了多少事”。推门时你沿运动方向使劲，门开了；拎着满满的水桶水平走路，胳膊再酸，水桶的高度一厘米也没涨。这两个例子提示了定义的形状：只有\emph{沿位移方向}的那份力才算数。于是我们定义：\emph{功}等于力的大小乘以物体沿力方向移动的距离。提桶走路时，提力向上、位移水平，两者垂直，这份力做的功是零——你的酸是肌肉内部的生理消耗，账并没有记到水桶头上。功像什么？像转账单：力是付款的动作，位移是付款的里程，乘起来是这一笔转了多少钱。它不像什么？它不像钱包本身——功不是物体“含有”的东西，物体不能说“我身上有一百焦的功”，就像你不能说“我身上有一百块钱转账”，转账只发生在过程中。

\textbf{第二个新概念：能量——账户余额。}有了转账单，就得有账户。我们定义：\emph{能量是一个系统“能够对外做功”的本领的度量}。被举高的锤子能砸桩，上紧的发条能走表，飞奔的子弹能穿板——它们账上都有余额。注意这个定义的妙处：它没有说能量“是什么东西”，只说它“能干什么”。能量不是一种看得见摸得着的物质，它是物理学家发明的一个\emph{记账量}；可千万别小看记账量，后面我们会看到，这是全宇宙管得最严的一本账。

\textbf{第三个新概念：动能——运动的余额。}子弹穿木板的麻烦现在可以解决了：运动物体账上的余额不跟速度成正比，而跟\emph{速度的平方}成正比。我们把这份余额叫动能，写成 \(\frac{1}{2}mv^2\)。速度翻倍，动能变成四倍——所以刹车距离变四倍，所以橡皮筋多拉两倍、弹出去的纸条获得四倍的动能（准确说是橡皮筋给出的能量变为四倍），飞得远得多。第二题的反直觉答案（b）就是这么来的。

\textbf{第四个新概念：势能——存起来的余额。}举高的锤子还没砸下去，能量在哪？既不在锤子的分子里，也不在空气里，它存在“锤子加地球”这个系统的\emph{位置关系}里。这种因位置而储存的能量叫势能。地面附近，重力势能就是重量乘高度；拉长的弹簧和橡皮筋存的是弹性势能，压缩或拉伸多少，存进去的账大致按形变量的平方记。势能像什么？像存在银行里的钱，取的时候连本带息（其实是分毫不差）能重新变回动能。它不像什么？它不像现金——你不能脱离“锤子—地球”或“弹簧—手”这个系统去谈势能，说“这个物体自己有重力势能”只是一种方便的简称，真正的账户开在系统上。

\textbf{第五个新概念：保守力与非保守力——可逆与不可逆的转账。}重力、弹簧力有个好脾气：你把物体举上去再原样放下来，它们做功一正一负，正好抵消，转出去的钱能全额转回来。这样的力叫\emph{保守力}，对它们才能定义势能。摩擦力正好相反：箱子推过去再推回来，摩擦力两次都倒着收“手续费”，来回一趟，账上的机械能平白少了一块。这样的力叫\emph{非保守力}。被摩擦收走的钱去哪了？没有消失——摸一摸急刹后的刹车盘、搓一搓冬天的双手，它们变热了。能量换了一种形态，存进了分子永不停歇的乱动里，这份账叫内能，那是热学篇章的主角。所以准确的说法是：\emph{机械能会亏，总能量从不亏}。摩擦消灭的不是能量，只是能量的“好用形态”。

\textbf{第六个新概念：功率——转账的速度。}同样爬三层楼，走上去和跑上去，转的账一样多（体重乘高度），可喘的程度天差地别。区别在于转账的\emph{快慢}：单位时间做的功叫功率。举重运动员举起杠铃只要一两秒，功率惊人；老太太用十分钟把同样重的东西挪上同样高的台阶，做的功完全一样，功率却小得多。电费单上的“度”是能量，电器铭牌上的“瓦”是功率——一个是账，一个是转账速度，别混。
\end{mainline}

六个概念齐了。回头看那本看不见的账：钢球被拉到鼻尖时，账上全是重力势能；荡到最低点，势能全部转成了动能；荡回另一头，动能再全额转回势能——这就是它认得原高度的原因。而每次少掉的那一点点高度，被空气阻力和转轴摩擦收进了内能账户，所以只低不高。转账可以，增发不行。

这套语言一旦上手，生活立刻变成一连串转账现场。拉弓：手臂肌肉的化学能变成弓臂的弹性势能，撒放瞬间又大部分变成箭的动能——古代强弓能把这个转换做得又快又狠，箭出弦的速度可以超过高铁。蹦床：人落下把动能压进床面的弹性势能，床面再把他弹回空中，每一跳亏掉的那部分变成了床布和空气的内能，所以不蹬就越来越低。钟摆、荡秋千、跳水板的起跳、撑杆跳的杆子，全是同一出戏：势能动能来回兑换，摩擦每笔抽佣。再看一个反着来的：自动扶梯匀速托着你上楼，你的动能没变，重力势能一直在涨——账是扶梯电动机以恒定的功率一笔笔转的；扶梯一旦停电，你就得用自己的化学能来付这笔钱，立刻感到腿沉。

\step{一句话模型}

\begin{oneline}
能量是宇宙唯一通行的货币：它不能被创造，也不能被销毁，只能在动能、势能、内能等账户之间转账——功和热是转账的两种方式，功率是转账的速度，而“只有保守力经手的账”才叫机械能守恒。
\end{oneline}

这句话里每一个词都对应上一步的一个定义。“货币”的比喻像什么？像货币之处在于总量守恒、形式可兑换、可储存可流通。它不像什么？它不像货币之处在于：能量没有通货膨胀，宇宙从不印钞；也没有任何银行能凭空放贷——这正是永动机不可能的原因。

\step{数学翻译器}

现在把上面的每一句白话翻译成式子。每条公式出来之前，先说它要回答什么；出来之后，立刻拿特殊情况敲打它。

\textbf{功：力的沿程积累。}要回答：一份恒力在一段直线路程上转了多少账？如果力 \(\mathbf F\) 与位移 \(\mathbf s\) 的夹角是 \(\theta\)，只有沿位移方向的分量 \(F\cos\theta\) 算数，于是
\[
W = Fs\cos\theta .
\]
检查特殊情况：\(\theta=0\)（顺着推），\(W=Fs\)，全额到账；\(\theta=90^{\circ}\)（提桶水平走路、圆周运动中绳的拉力），\(\cos\theta=0\)，\(W=0\)——向心方向的力只拐弯不转账，这是圆周运动速率不变（无其他切向力时）的能量解释；\(\theta=180^{\circ}\)（滑动摩擦力），\(W=-Fs\)，负号表示账从物体的机械能里被划走。功是标量，但有正负，正负就是“转入还是转出”。

\begin{center}
\begin{tikzpicture}[scale=1.0,>={Stealth[length=3mm]}]
  \draw[fill=gray!15] (0,0) rectangle (1.2,0.8);
  \draw[->,very thick] (0.6,0.4) -- (3.1,1.84) node[above left]{$\mathbf F$};
  \draw[->,very thick] (0.6,0.4) -- (3.6,0.4) node[midway,below]{$\mathbf s$};
  \draw[dashed] (2.6,0.4) arc[start angle=0,end angle=30,radius=2.0];
  \node at (2.45,0.75) {$\theta$};
  \draw[dashed] (3.1,1.84) -- (3.1,0.4);
  \draw[->,thick] (0.6,0.4) -- (3.1,0.4) node[midway,above=2pt]{\small $F\cos\theta$ 有效的部分};
\end{tikzpicture}
\end{center}

\textbf{动能：为什么是速度的平方。}要回答：运动的物体账上有多少余额？用我们已有的知识就能推出来。设恒力 \(F\) 推着质量 \(m\) 的物体从静止出发，沿力的方向移动距离 \(s\)。牛顿第二定律给出加速度 \(a=F/m\)，匀加速运动学给出 \(v^2=2as\)。两式联立消去 \(a\)：
\[
Fs=\frac{1}{2}mv^2 .
\]
左边正是外力转进来的账，右边就只能是物体账上新增的余额。于是动能
\[
E_{\mathrm k}=\frac{1}{2}mv^2 .
\]
检查：\(v=0\) 时动能为零，合理；\(v\) 变成两倍，动能变四倍——速度 \(80\) 千米每小时的刹车距离约为 \(40\) 千米每小时的四倍，不是两倍。这就是为什么高速公路上“保持车距”的铁律要按速度的平方来留余量，第二题的反直觉就此翻案。

\begin{mathbridge}
\textbf{变力的功：面积与积分。}力随位置变化时，把路程切成无数小段，每段里力几乎不变，功是 \(\Delta W\approx F(x)\Delta x\)，全程的功就是 \(F\)-\(x\) 图线下方的面积：
\[
W=\int_{x_1}^{x_2}F(x)\,\mathrm dx .
\]
用同样的“无限切分再求和”处理变加速情形，可以证明更一般的\emph{动能定理}：合外力对物体做的总功等于物体动能的增量，
\[
W_{\text{合}}=\Delta E_{\mathrm k}=\frac{1}{2}mv_2^2-\frac{1}{2}mv_1^2 .
\]
注意这条定理对恒力、变力、直路、弯路一律成立——它正是“能量记账法”比“逐瞬间受力分析”省事的根源：只问首尾，不问过程。
\end{mathbridge}

\textbf{重力势能：重量乘高度。}要回答：把物体举高 \(h\)，账上存进多少？克服重力做功 \(mgh\)，全数存入，故
\[
E_{\mathrm p}=mgh ,
\]
其中 \(h\) 从任意选定的零高度量起。检查：\(h=0\) 处势能为零——零点开在哪纯属记账习惯，只有高度\emph{差}有物理意义；\(h\) 取负也合法，地下室里的物体势能就是负的，账上是“欠”的，落差照样能取出来用。水电站利用的正是落差：水量乘落差乘 \(g\)，就是水库账上的总余额。

\textbf{弹性势能：面积里的二分之一。}要回答：把弹簧从原长拉伸 \(x\)，存进多少？胡克定律说拉力随形变线性增长：\(F=kx\)。力不是恒定的，从一端的 \(0\) 涨到另一端的 \(kx\)，平均值是 \(\frac{1}{2}kx\)，乘以距离 \(x\)：
\[
E_{\mathrm p}=\frac{1}{2}kx^2 .
\]
这个二分之一不是凑出来的，它就是力—位移图里那个三角形的面积。检查：\(x=0\) 时为零，合理；\(\pm x\) 给出同样的值——压缩和拉伸存一样多的账，橡皮筋弹射实验里“拉长两倍、能量四倍”的谜底就在这个平方里。

\begin{center}
\begin{tikzpicture}[scale=0.9,>={Stealth[length=2.5mm]}]
  \draw[->] (0,0) -- (4.6,0) node[right]{$x$};
  \draw[->] (0,0) -- (0,3.4) node[above]{$F$};
  \draw[very thick] (0,0) -- (3.6,2.7) node[right]{$F=kx$};
  \draw[dashed] (3.6,0) node[below]{$x$} -- (3.6,2.7) -- (0,2.7) node[left]{$kx$};
  \fill[pattern=north east lines,pattern color=blue!60] (0,0) -- (3.6,2.7) -- (3.6,0) -- cycle;
  \node at (2.55,0.85) {$E_{\mathrm p}=\dfrac12 kx^2$};
\end{tikzpicture}
\end{center}

\textbf{机械能守恒：只在保守力经手时成立。}要回答：什么条件下动能加势能的总额不变？把只有保守力做功的过程逐笔记账——重力转的账：动能涨多少、势能就跌多少——可得
\[
E_{\mathrm k}+E_{\mathrm p}=\text{常量}\qquad(\text{只有保守力做功}) .
\]
检查极端情形：摆的最高点，\(v=0\)，账上全是势能；最低点，势能最少、动能最大；两点之间，两边此消彼长、总额不动。一旦有摩擦、空气阻力经手，等式右边就不再是常量，而是逐笔记亏——亏掉的那部分没有蒸发，它进了内能账户。这才是完整的守恒律：\emph{把内能也算进来，总账永远平衡}。

\textbf{功率：转账的快慢。}要回答：一段时间里账转得多快？
\[
P=\frac{W}{t},\qquad\text{恒力与速度同向时还可写成 }P=Fv .
\]
检查 \(P=Fv\) 的两个极端：推一堵纹丝不动的墙，力再大，\(v=0\)，功率为零——墙账上一焦耳也没收到；起重机吊着集装箱\emph{匀速}上升，合力为零但钢索的拉力不为零，拉力以 \(P=Fv\) 的速率持续向集装箱—地球系统转账。后一式还解释了一个驾驶常识：车提速到两倍，发动机维持同样牵引力所需的功率也是两倍，而克服空气阻力所需的功率涨得还要更快——这是高速费油的第一层原因。

\begin{mathbridge}
\textbf{功率的微商形式。}把“快慢”说得更精确：功率是功对时间的变化率，
\[
P=\frac{\mathrm dW}{\mathrm dt}=\mathbf F\cdot\frac{\mathrm d\mathbf s}{\mathrm dt}=\mathbf F\cdot\mathbf v ,
\]
点乘自动挑出“沿运动方向”的力分量。于是前面所有关于夹角的讨论被这一个符号统一收编：垂直不做功，反向做负功。
\end{mathbridge}

\textbf{电费单上的账。}家家都交的电费，是能量记账法最接地气的应用。功率 \(1\) 千瓦的电器开 \(1\) 小时，消耗 \(1\) 千瓦时，俗称一度电：
\[
1\ \text{度}=1000\ \text{瓦}\times3600\ \text{秒}=3.6\times10^{6}\ \text{焦} .
\]
这笔钱有多大？全部拿来克服重力做功，能把一吨水垂直抬升约 \(360\) 米。一只 \(1500\) 瓦的电热水壶烧开一壶水大约 \(5\) 分钟，用掉约 \(0.12\) 度；一台变频空调制冷季一天几度的电费，本质是压缩机整晚以几百瓦的功率持续把热量从屋里“搬”到屋外——注意它付的是搬运费，不是热量的价钱，这正是空调比电暖器省电的深层原因，热学篇章会细说。

\textbf{一次带真实数据的估算。}一碗米饭的化学能大约是 \(1\times10^{6}\) 焦（约 \(250\) 千卡）。假如你的身体能把这些能量百分之百转成爬楼的机械功，一个 \(60\) 千克的人能爬多高？由 \(mgh=10^{6}\) 焦，取 \(g\approx10\) 米每二次方秒，得
\[
h=\frac{10^{6}}{60\times10}\approx 1700\ \text{米} ,
\]
相当于两座泰山的海拔爬升。当然身体远不是百分之百的效率，大部分能量最终变成了体温和汗液蒸发的内能——但数量级是诚实的：食物里的化学账，远比我们感觉的“一碗饭”厚得多。再看功率：同一个人 \(15\) 秒爬上 \(10\) 米高的楼梯，机械功率约 \(60\times10\times10/15=400\) 瓦，短时间输出能超过半台微波炉的额定功率；而持续一整天，人的平均功率只剩一百瓦上下——一匹真正的马能持续输出七百多瓦，“马力”这个单位就是这么来的。

\textbf{一项工程应用：把刹车变成充电。}传统汽车刹车，是把整车动能通过刹车片摩擦全数转成内能，白白烧掉。电动车和高铁的再生制动反其道而行：刹车时让车轮倒拖电动机，电动机临时改当发电机，把动能转回电池或电网。算一笔账：一辆 \(1.5\) 吨的车以 \(100\) 千米每小时（约 \(28\) 米每秒）行驶，动能为
\[
E_{\mathrm k}=\frac12\times1500\times28^{2}\approx 5.9\times10^{5}\ \text{焦}\approx 0.16\ \text{度} .
\]
单次刹车似乎不多，可市区通勤一天几十上百次启停，回收的电量足以让续航多出可观的一截；高铁一次进站制动回收的能量，够一列车的照明和空调用上一阵。这不是创造能量，只是拒绝把账烧掉——守恒律没破，浪费少了。

\step{模型的边界}

每条守恒律都是一份合同，合同都有条款。能量记账法的边界主要有五条。

第一，\textbf{机械能守恒有条件}：只有保守力（重力、弹力、万有引力这类）做功时，动能加势能才是常量。有摩擦、有碰撞形变、有人推发动机转，等式就要改写成“机械能的变化等于非保守力做的功”。把过山车全程当机械能守恒来设计，会低估第一站的高度、酿成大祸——工程师必须把每一段摩擦损耗都估进账里。

第二，\textbf{总能量守恒要求系统封闭}。说“能量守恒”之前先问：账本的边界画在哪？只画小车，摩擦把账转给了地面和空气，小车的账当然亏；把地面、空气、声音、热量全圈进来，总账才平。永动机的第一种死法正在于此：设计者声称机器对外发电还能自己转下去，等于宣称一个封闭账本上的总额在涨——两百年里各国专利局收到的这类方案数以千计，没有一个通过实测。法国科学院早在 1775 年就宣布不再审理永动机专利，不是傲慢，是这本账从来没有出过一分钱的错。要注意，还有第二种更隐蔽的永动机：不违反能量守恒，只想把海水里的内能无条件地全部取来做功——它的死刑要等到热学篇章里熵出场时才正式宣判，这里先留个悬念。

第三，\textbf{势能的零点和归属}。势能零点任意，只有差值有物理意义；势能属于系统而不属于单个物体。说“书的重力势能是多少焦”必须默认了“以桌面为零点、在书—地球系统里”这两层省略，否则就是病句。

第四，\textbf{几句日常语言的纠正}。“电池没电了”的真实含义不是能量被用光了，而是电池里可逆转成电能的化学能余额见底，转出去的账一分不少地变成了灯光、声波和手机微微的发热。“这个机器消耗能量”也是同一句口语的变体——能量从不被消耗，只是被转走，而且常常转成了我们不想要的形态。反过来，“精力充沛”“元气满满”这类说法里的“能量”是心理隐喻，跟物理学那个可测量、可转账、守恒的记账量不是一回事，两者互不背书。还要警惕把能量当成某种神秘流体的想法：能量不是装在物体里的“东西”，它是我们为一切变化统一记账而发明的量；它的“实在性”体现在账永远平，而不是体现在你能舀一勺能量出来看看。

第五，\textbf{经典记账法的适用范围}。本章全部公式默认速度远小于光速。接近光速时，\(\frac{1}{2}mv^2\) 不再够用，能量与质量的账本要按相对论重记——挑战层盒子里先剧透一行。微观世界里能量依然守恒，只是账户变成了分立的“能级”，那是量子篇章的故事。还有一点值得强调：能量守恒在本书不是一条“恰好没发现反例”的经验规律。 二十世纪初，数学家诺特证明了一个深刻的定理：只要物理规律不随时间平移而改变——今天做实验和明天做实验结果一样——就必然存在一个守恒量，它正是能量。能量守恒是“时间均匀性”的收据。

\begin{challenge}
\textbf{动能平方律的更深来历。}相对论的记账方式是：一个物体的总能量写成 \(E=\gamma mc^{2}\)，其中 \(c\) 是光速，\(\gamma=1/\sqrt{1-v^{2}/c^{2}}\) 在低速时略大于 \(1\)。把 \(\gamma\) 按 \(v^{2}/c^{2}\) 的幂次展开，第一项是常数 \(mc^{2}\)（物体静止时账上就有的巨额余额），第二项正好是 \(\frac{1}{2}mv^{2}\)——本章的动能公式原来是相对论账本在低速下的第一条修正。更高阶的项在粒子加速器里再也不是小字：把质子加速到光速的 \(99.9999991\%\)，它的动能是经典公式预言值的几千倍。账本的框架没有变，变的是每种账户的换算汇率。
\end{challenge}

\step{回头解释开场现象}

现在回到教室里的那只钢球，把本章走完最后一环。

松手瞬间，钢球在鼻尖高度，速度为零：账上全部是重力势能 \(mgh\)。下摆过程中重力持续做正功，势能逐笔记成动能；最低点时高度最小、速度最大，账上几乎全是动能。上摆过程反过来，动能逐笔兑回势能。如果世界是理想封闭的——没有空气阻力，吊点没有摩擦——钢球将精确地回到 \(h\)，一毫米不多、一毫米不少。\emph{它“记得”高度，其实是机械能守恒在替它记账}：任何时刻，动能加势能等于最初的 \(mgh\)，而回到鼻尖那一侧的最高点时动能必须为零（速度在那里反向），于是势能必须全额复原，高度也就必须复原。

现实里空气阻力和转轴摩擦每笔都抽走一点手续费，转成钢球、空气和轴承的内能，所以钢球每次回来都略低一点。也正因为只能低、不能高，教授才敢拿脸做抵押——如果有什么机制能让摆自己越荡越高，那等于账本会自己涨钱，整个物理学的大楼就得塌掉重盖。

那秋千为什么又能越荡越高？因为荡秋千的人不是封闭账本。你在最高点下蹲、在最低点站起，蹬腿和摆身体的每个动作都在用肌肉的化学能对秋千—人系统做正功，而且恰好踩在摆动节奏最敏感的时间点上——每次转进的账都比摩擦抽走的多，振幅就越攒越大。一旦人停下动作，账本上只剩摩擦在收手续费，秋千便无可挽回地低下去。“在对的时间转账”这件事本身就是一门大学问，它叫共振，振动篇章里会专门讲；这里只需要记住：摆幅的涨落从来不是账本的奇迹，只是有人在持续充值。

第一题也一并了结：无摩擦时，直坡和弯道底部的高度差相同，重力转出的账相同，由 \(\frac{1}{2}mv^2=mgh\) 得两球速率\emph{完全一样}，与路径形状无关——答案（c）。这就是保守力的脾气：只认首尾，不认路。（顺带一提：若把摩擦放回来，弯道因为路程更长、挤压更狠，底部反而更慢。直觉又一次在“过程不重要”这件事上栽了跟头。）至于钥匙摆被笔杆拦住后仍爬回原高度：半径变了、轨迹变了、最低点速度没变、高度差没变——账本只看高度差。

\step{脑洞实验与挑战题}

\begin{thought}[从月亮上扔一块石头]
脑洞实验：想象你站在月球轨道上（距地心约 \(3.8\times10^{8}\) 米），轻轻\emph{丢下}一块石头——不推，只是松手，并且先忽略大气。石头在引力下一路加速坠落地球，撞地前一刻速度大约多少？用记账法而不必逐段追踪：引力势能从“远”到“近”减少，全部转成动能。完整的引力势能公式（第 8 章的礼物）给出的答案是约 \(11\) 千米每秒——恰好是第二宇宙速度。这不是巧合：从地面以 \(11\) 千米每秒竖直上抛的物体恰好能挣脱地球飞到无穷远，而“掉进来”正是“飞出去”的时间倒放，两笔账必然相等。再推一步：这块石头进入大气层后，那笔巨额动能几乎全数转成内能——这就是流星为什么会烧起来，也是航天器返回舱需要防热大底的全部原因。一次“轻轻松手”，落地时却是一团火球，能量账本对浪漫的想象力毫不留情。
\end{thought}

\begin{puzzles}
\item \textbf{过山车圆环。}过山车的圆环半径为 \(R\)，小车从一段光滑高坡上由静止释放，要安全通过圆环最高点而不脱轨，释放点至少要比圆环最高点高多少？（提示：最高点的“不脱轨”条件不是速度大于零，而是轨道还得提供向心作用——临界情形是轨道对小车的支持力恰好为零，重力独自提供所需的向心加速度 \(v^{2}/R\geq g\)。把这条动力学条件和机械能守恒联立，答案是再高出 \(R/2\)。真实的过山车要留安全余量，所以释放坡总是明显高于这个临界值。）
\item \textbf{两支箭。}同一位置、同一速率，一支箭竖直向上射，一支箭水平射出，不计空气阻力。落地瞬间，哪支箭的速率更大？（提示：别去算轨迹。问账本：两次的起始动能相同，落回地面的高度差相同，重力转的账相同。竖直上抛那支箭在最高点停了一瞬间、绕了远路——可保守力只认首尾。）
\item \textbf{子弹与半块木板。}一颗子弹恰好能击穿一块木板，穿出时速度降为零。换一块厚度一半的同款木板，同一颗子弹以同样的速度射入，穿出时还剩多少速度？（提示：木板厚度减半，摩擦力做功减半，即动能只损失一半。剩下的是一半动能，不是一半速度——速度剩下 \(1/\sqrt{2}\)，约七成。这就是正文“麻烦二”的答案。）
\item \textbf{抽水蓄能的账。}抽水蓄能电站深夜用富余的电把下水库的水抽进上水库，白天用电高峰再放水发电。一座电站上水库比下水库高约 \(500\) 米，库存水 \(10^{10}\) 千克。全部放完，理论上能发多少电（用“度”，即千瓦时作单位）？再想一想：这套“抽上去再放下来”的循环明明每一步都有损耗，为什么电网还愿意修它？（提示：\(E=mgh=10^{10}\times10\times500=5\times10^{13}\) 焦，除以每度 \(3.6\times10^{6}\) 焦，约一千四百万度。损耗换到的是时间——把半夜没人要的电存成白天抢手的电，能量打八折，价值翻几倍。守恒的账和经济的账，是两种不同的记账法。）
\end{puzzles}

下一章我们要遇到另一本账：碰撞发生时，能量这本账经常显得“对不平”——两团橡皮泥撞在一起，动能凭空少了一大半，可有一个量却纹丝不动。那个量叫动量。能量和动量，是宇宙记账本上的两页，谁也不能替代谁。
